\subsection{推广到多模复形；典范同构}

设 B, B' 为两个环，C 为 (A, B)-双模复形，C' 为 (A, B')-双模复形；则 (Homgr\_A (C, C′), D) 是 (B, B′)-双模复形，且典范同态 \(\lambda: H(\text{Homgr}_A (\text{C}, \text{C'})) \to \text{Homgr}_A (\text{H}(\text{C}), \text{H}(\text{C'}))\) 是 (B, B')-双模同态。

若 M 是 (A, B)-双模，N 是 (A, B')-双模，则 Homgr \(_{A}\) (L(M), I(N)) 是 (B, B')-双模复形，因此 Ext \(_{A}\) (M, N) 具有自然的分次 (B, B')-双模结构；在次数 0 项上，该结构与 Hom \(_{A}\) (M, N) 的结构一致（II，第 35 页）。

若 \(\lambda\in\mathbf{B},\lambda^{\prime}\in\mathbf{B}^{\prime}\) 并且如果我们记 \(\lambda_{\mathrm{M}},\lambda_{\mathrm{N}}^{\prime},\lambda_{\mathrm{E}},\lambda_{\mathrm{E}}^{\prime}\) 为位似 \(x\mapsto x\lambda,y\mapsto y\lambda^{\prime},z\mapsto\lambda z,z\mapsto z\lambda^{\prime}\) ，它们分别是 M, N, Ext \(_{\mathrm{A}}\) (M, N), Ext \(_{\mathrm{A}}\) (M, N) 的，则我们有

\[
\lambda_ {\mathrm{E}} = \operatorname{Ext} _ {\mathbf {A}} \left(\lambda_ {\mathbf {M}}, 1\right), \quad \lambda_ {\mathrm{E}} ^ {\prime} = \operatorname{Ext} \left(1, \lambda_ {\mathbf {N}} ^ {\prime}\right),
\]

这给出了 Ext \(_{A}\) (M, N) 的双模结构的另一种描述。我们把第 4 号和第 6 号推广到多模复形的情形留给读者。

设 C, C′ 和 C′′ 是 A-模的复形。映射的复合定义了一个次数为零的分次同态：

\[
\operatorname{Homgr} _ {\mathbf {A}} \left(\mathrm{C} ^ {\prime}, \mathrm{C} ^ {\prime \prime}\right) \otimes_ {k} \operatorname{Homgr} _ {\mathbf {A}} \left(\mathrm{C}, \mathrm{C} ^ {\prime}\right)\rightarrow \operatorname{Homgr} _ {\mathbf {A}} \left(\mathrm{C}, \mathrm{C} ^ {\prime \prime}\right).\tag{14}
\]

设 B、B'、E 为环，C、C'、C'' 分别为 (B', A)-双模、(A, E)-双模、(B, E)-双模的复形。通过限制 II 第 73 页的典范同构，得到 (B, B')-双模的一个双射同态：

\[
\operatorname{Homgr} _ {E} \left(C \otimes_ {A} C ^ {\prime}, C ^ {\prime \prime}\right)\rightarrow \operatorname{Homgr} _ {A} \left(C, \operatorname{Homgr} _ {E} \left(C ^ {\prime}, C ^ {\prime \prime}\right)\right).\tag{15}
\]

最后，设 B 为一个环，C 为右 B-模的复形，C' 为右 A-模的复形，C'' 为 (B, A)-双模的复形。我们由典范同态（II，第 75 页）推出

\[
\mathrm{C} _ {p} \otimes_ {\mathrm{B}} \operatorname{Hom} _ {\mathrm{A}} \left(\mathrm{C} _ {q} ^ {\prime}, \mathrm{C} _ {r} ^ {\prime \prime}\right)\rightarrow \operatorname{Hom} _ {\mathrm{A}} \left(\mathrm{C} _ {q} ^ {\prime}, \mathrm{C} _ {p} \otimes_ {\mathrm{B}} \mathrm{C} _ {r} ^ {\prime \prime}\right)
\]

一个次数为零的分次同态：

\[
\mathrm{C} \otimes_ {\mathrm{B}} \operatorname{Homgr} _ {\mathrm{A}} \left(\mathrm{C} ^ {\prime}, \mathrm{C} ^ {\prime \prime}\right)\rightarrow \operatorname{Homgr} _ {\mathrm{A}} \left(\mathrm{C} ^ {\prime}, \mathrm{C} \otimes_ {\mathrm{B}} \mathrm{C} ^ {\prime \prime}\right).\tag{16}
\]

当 C 是有限生成投射模时，此同态是双射的（II，第 75 页，命题 2）。

命题12.——同态（14）、（15）、（16）是复形的态射。

我们来证明它，例如对于同态（14）。记

\[
\kappa : \operatorname{Homgr} _ {\mathbf {A}} \left(\mathrm{C} ^ {\prime}, \mathrm{C} ^ {\prime \prime}\right) \otimes_ {k} \operatorname{Homgr} _ {\mathbf {A}} \left(\mathrm{C}, \mathrm{C} ^ {\prime}\right)\rightarrow \operatorname{Homgr} _ {\mathbf {A}} \left(\mathrm{C}, \mathrm{C} ^ {\prime \prime}\right)
\]

这个同态。设 \(f \in \operatorname{Homgr}_{\mathbf{A}}(\mathbf{C}', \mathbf{C}^{\prime})_{p}\) 和 \(g \in \operatorname{Homgr}_{\mathbf{A}}(\mathbf{C}, \mathbf{C}^{\prime})_{q}\) ；于是根据定义我们有 \(\kappa(f \otimes g) = f \circ g\) 。此外：

\[
\begin{array}{r l} \mathrm{D} (f \otimes g) = & \mathrm{D} f \otimes g + (- 1) ^ {p} f \otimes \mathrm{D} g = (d ^ {\prime \prime} \circ f) \otimes g - (- 1) ^ {p} (f \circ d ^ {\prime}) \otimes g + \\ & + (- 1) ^ {p} f \otimes (d ^ {\prime} \circ g) - (- 1) ^ {p + q} f \otimes (g \circ d), \end{array}
\]

由此可得

\[
\begin{array}{r l} & {\kappa \big (\mathbf {D} (f \otimes g) \big) = d ^ {\prime \prime} \circ f \circ g - (- 1) ^ {p} f \circ d ^ {\prime} \circ g +} \\ & {\qquad + (- 1) ^ {p} f \circ d ^ {\prime} \circ g - (- 1) ^ {p + q} f \circ g \circ d} \\ & {\qquad = d ^ {\prime \prime} \circ f \circ g - (- 1) ^ {p + q} f \circ g \circ d = \mathbf {D} (f \circ g) = \mathbf {D} \big (\kappa (f \otimes g) \big).} \end{array}
\]

类似地可证明同态(15)和(16)是复形的态射。

由态射 (14) 我们推出 \(k\) -模的同态（X，第 62 页）

\[
\mathrm{H} ^ {p} \left(\operatorname{Homgr} _ {\mathbf {A}} \left(\mathrm{C} ^ {\prime}, \mathrm{C} ^ {\prime \prime}\right)\right) \otimes_ {k} \mathrm{H} ^ {q} \left(\operatorname{Homgr} _ {\mathbf {A}} \left(\mathrm{C}, \mathrm{C} ^ {\prime}\right)\right)\rightarrow \mathrm{H} ^ {p + q} \left(\operatorname{Homgr} _ {\mathbf {A}} \left(\mathrm{C}, \mathrm{C} ^ {\prime \prime}\right)\right). \tag {17}
\]

取 C = A，我们看到同态：

\[
\operatorname{Homgr} _ {\mathbf {A}} \left(\mathrm{C} ^ {\prime}, \mathrm{C} ^ {\prime \prime}\right) \otimes_ {k} \mathrm{C} ^ {\prime} \rightarrow \mathrm{C} ^ {\prime \prime}\tag{18}
\]

它把 \(f \otimes x\) 映到 \(f(x)\) ，是左 A-模复形的一个态射；与之相伴的是分次 A-模的一个典范同态（X，第 80 页） \(\gamma : \mathrm{H}(\mathrm{Homgr}_{\mathbf{A}}(\mathbf{C}', \mathbf{C}')) \otimes_k \mathrm{H}(\mathbf{C}') \to \mathrm{H}(\mathbf{C}'')\) ，它对应于 \(k\) -模的典范同态

\[
\lambda : \mathrm{H} (\operatorname{Homgr} _ {\mathbf {A}} \left(\mathrm{C} ^ {\prime}, \mathrm{C} ^ {\prime \prime}\right)) \rightarrow \operatorname{Homgr} _ {\mathbf {A}} \left(\mathrm{H} (\mathrm{C} ^ {\prime}), \mathrm{H} (\mathrm{C} ^ {\prime \prime})\right).
\]

